\subsection{数学实例化模板：哪些对象可承载物理型时空骨架}\label{subsec:typed-address-biaxial-completion-instantiation-template}
上一节已经把粗糙时空地形图压缩为四元组
\[
\mathfrak G=(\preceq,\tau,d_{\mathrm{res}},\Omega),
\]
并说明这并不是现成的 Minkowski、Lorentz 或 Einstein 几何，而只是统一上位逻辑在观察者化、局部化、因果化与资源化之后得到的最小前几何骨架。
由此随之而来的问题是：究竟哪些具体数学对象有资格承载这样一套骨架？

本文第 \(15\) 章的边界已经很清楚：当前稿件并不直接推出 Minkowski 度规、Lorentz 对称、Einstein 型场方程或测量公设。
这些都属于更强的实例化层结论，只有当某个具体数学对象同时保持 forcing 结构、局部性、因果相容、失败模式、多尺度一致性、谱接口与资源代价时，才可能由当前骨架进一步收缩出来。
换句话说，并不是任何数学对象都配得上被称为“物理型实例化”；它必须满足一组比“能够编码数据”更强的结构条件。

本节的任务，就是把这一实例化层要求写成一个明确模板。
其目的有二：第一，给出一条进入物理解释之前的准入条件，避免把纯形式对象误说成时空模型；第二，把此前分散出现的观察者状态、restriction 逆系统、局部截面、谱接口、残差外置与资源预算统一压成一组可直接调用的实例化数据。

\subsubsection{结构保真先于对象代入}\label{subsubsec:typed-address-biaxial-completion-instantiation-structure-faithfulness}
从逻辑上说，上一节的四元组 \(\mathfrak G\) 并不会从任意集合、任意概率空间或任意流形中自动跳出来。
若要真正抽取出 \((\preceq,\tau,d_{\mathrm{res}},\Omega)\)，至少要有六类对象可用：
\begin{enumerate}
  \item 状态、forcing、更新与通信结构，否则无法定义可达/因果前序；
  \item 局部读出、事件定位与共同支撑，否则无法定义时间投影与局部比较；
  \item 多尺度一致性，否则无法把不同分辨率解释为同一本体的不同精细度；
  \item 局部截面与重叠转移，否则无法谈全局化与 Čech 障碍；
  \item 残差外置与失败证书，否则无法区分“对象不存在”与“协议失败”；
  \item 谱接口与资源预算，否则无法把高阶统计、同步成本与账本代价组织成几何量。
\end{enumerate}
因此，实例化并不是“随便找一个对象去附会物理”，而是寻找一个具体对象族，使它能够逐项承载这六类结构。
只有在这一意义下，上一节的前几何四元组才从抽象记号下降为真正的对象层几何读法。

\subsubsection{物理型实例化数据}\label{subsubsec:typed-address-biaxial-completion-physical-instantiation-data}
为了把这六类结构写成单一模板，先引入统一定义。

\begin{definition}[物理型实例化数据]\label{def:typed-address-biaxial-completion-physical-instantiation-data}
称六元组
\[
\mathcal I=
(\mathcal P,\mathcal X,\mathcal U,\mathcal R,\mathcal E,\mathcal C)
\]
为一组物理型实例化数据，若其中：
\begin{enumerate}
  \item \(\mathcal P\) 是状态—forcing—更新层，至少包括
  \[
  (P,\le,\Vdash,\{\mathrm{Upd}_\alpha\},\{\mathrm{Com}_\mu\},\iota,\lambda,\varrho),
  \]
  其中 \(P\) 为状态空间，\(\le\) 为精化或信息序，\(\Vdash\) 为 forcing 关系，\(\mathrm{Upd}_\alpha\) 为内部更新，\(\mathrm{Com}_\mu\) 为通信/传输，\(\iota\) 为观察者索引，\(\lambda\) 为事件定位，\(\varrho\) 为局部支撑；
  \item \(\mathcal X\) 是多尺度读出层，至少包括一个逆系统
  \[
  \cdots \longrightarrow X_{m+1}\xrightarrow{\pi_{m+1,m}}X_m\longrightarrow\cdots\longrightarrow X_1,
  \]
  及其逆极限对象 \(X_\infty\)，并带有与状态层相容的读出映射；
  \item \(\mathcal U\) 是局部覆盖—截面层，至少包括覆盖 \(\{U_i\}\)、局部可读对象的预层/层 \(\mathscr F\)、局部截面与重叠上的转移数据；
  \item \(\mathcal R\) 是谱—统计层，至少包括由可见数据抽取的高阶矩、压力、primitive 轨道谱、\(\zeta\) 接口或等价谱对象；
  \item \(\mathcal E\) 是残差—外置层，至少包括纤维残差、\(\mathrm{NULL}\)/失败证书，以及必要时的 prime-register 或其他可逆账本；
  \item \(\mathcal C\) 是资源—代价层，至少包括局部比较所需的预算向量与把预算标量化的规则。
\end{enumerate}
\end{definition}

定义 \ref{def:typed-address-biaxial-completion-physical-instantiation-data} 的六个分量分别承担不同职责。
\(\mathcal P\) 决定“谁在更新、谁能比较、谁能通信”；
\(\mathcal X\) 决定“同一对象在不同分辨率下如何相容”；
\(\mathcal U\) 决定“局部可读为何未必能全局拼起”；
\(\mathcal R\) 决定“哪些高阶统计或 primitive 指纹是机制不变量”；
\(\mathcal E\) 决定“哪些自由度被折叠抹去、却仍必须外置承载”；
\(\mathcal C\) 则把这一切变成真正的几何预算。

\subsubsection{允许的实例化条件}\label{subsubsec:typed-address-biaxial-completion-admissible-instantiation}
仅有六元组本身还不够；它还必须满足若干兼容性条件，才能被视为本文意义下可进入物理解释的实例。

\begin{definition}[允许的物理型实例化]\label{def:typed-address-biaxial-completion-admissible-instantiation}
称物理型实例化数据 \(\mathcal I\) 为允许的，若它满足下列六条：
\begin{enumerate}
  \item \textbf{观察者对称与无特权条件：}
  观察者标签的重命名可被整体结构自同构吸收，因而观察者只是状态纤维的索引，而不是形而上的特权主体；
  \item \textbf{因果相容更新条件：}
  状态更新与通信共同定义一个可达前序，并与事件定位相容，不产生回溯性矛盾；
  \item \textbf{多尺度一致性条件：}
  多尺度读出满足 restriction 逆系统的一致性，并具有逆极限对象；
  \item \textbf{局部—全局障碍条件：}
  局部截面存在并不自动推出全局对象存在，且重叠失败可被显式记录成 Čech 型残差、holonomy 或等价粘接障碍；
  \item \textbf{残差外置条件：}
  任何因折叠、压缩或局部读出而被抹掉、但在可逆恢复中不可忽略的自由度，都必须在 \(\mathcal E\) 中被外置承载，而不能被“免费忽略”；
  \item \textbf{谱—资源可审计条件：}
  从可见层抽取出的高阶矩、压力、primitive 轨道谱与资源预算之间必须存在可核验接口，而不是事后经验拟合。
\end{enumerate}
\end{definition}

定义 \ref{def:typed-address-biaxial-completion-admissible-instantiation} 可被视为实例化的准入门槛。
若某个对象只满足其中一两条，例如它只会产生一组概率分布、只带一个预设流形结构，或只提供某种实现层编码，那么它仍然可以是有趣模型，但还不够成为本文意义下的物理型实例化。

\subsubsection{从实例化数据抽取粗糙时空四元组}\label{subsubsec:typed-address-biaxial-completion-instantiation-to-rough-spacetime}
一旦 \(\mathcal I\) 满足定义 \ref{def:typed-address-biaxial-completion-admissible-instantiation}，上一节给出的粗糙时空四元组就不再是悬空记号，而可以系统地从 \(\mathcal I\) 中抽出来。

\begin{proposition}[实例化到粗糙时空骨架的抽取]\label{prop:typed-address-biaxial-completion-instantiation-extraction}
\leanverified{paper\_typed\_address\_biaxial\_completion\_instantiation\_extraction}
设 \(\mathcal I\) 是一组允许的物理型实例化。
则它可自然诱导一个粗糙时空四元组
\[
\mathfrak G(\mathcal I)
=
(\preceq_{\mathcal I},\tau_{\mathcal I},d_{\mathrm{res},\mathcal I},\Omega_{\mathcal I}),
\]
其中：
\begin{enumerate}
  \item \(\preceq_{\mathcal I}\) 由更新与通信结构给出可达/因果前序；
  \item \(\tau_{\mathcal I}\) 由事件定位映射 \(\lambda\) 及其时间投影给出；
  \item \(d_{\mathrm{res},\mathcal I}\) 由资源预算 \(\mathcal C\) 的标量化产生；
  \item \(\Omega_{\mathcal I}\) 由局部截面重叠、闭路 residual 与跨尺度不交换共同产生。
\end{enumerate}
\end{proposition}

\begin{proof}
定义 \ref{def:logic-expansion-observer-spacetime-state-model} 已经给出状态层中的 \((P,\iota,\lambda,\varrho)\) 与时间投影 \(\tau=\chi\circ\lambda\)；
定义 \ref{def:logic-expansion-causal-compatibility} 则给出更新与通信如何与因果顺序相容。
因此 \(\preceq_{\mathcal I}\) 可取为“由有限次允许更新与通信达到”的可达闭包，\(\tau_{\mathcal I}\) 则由 \(\lambda\) 的时间分量直接给出。

资源准距离部分由上一节的定义 \ref{def:typed-address-biaxial-completion-resource-vector} 与定义 \ref{def:typed-address-biaxial-completion-resource-quasidistance} 直接抽取。
而 \(\Omega_{\mathcal I}\) 的覆盖方向来自定义 \ref{def:prefix-site-h2-gluing-obstruction} 及其后的局部截面—Čech 障碍链；尺度方向则来自 restriction 逆系统上的跨层不交换 residual。
因此，\(\mathcal I\) 一旦满足上述兼容性，\(\mathfrak G(\mathcal I)\) 便可自然定义。证毕。
\end{proof}

命题 \ref{prop:typed-address-biaxial-completion-instantiation-extraction} 说明：上一节的粗糙时空骨架并不是脱离对象层的抽象比喻，而是每一组合格实例化数据都自动带出来的最小几何读法。

\subsubsection{五类自然实例簇}\label{subsubsec:typed-address-biaxial-completion-five-instantiation-clusters}
在这一模板下，最自然的实例并不是唯一的，而是分成若干“实例簇”；它们各自擅长承载不同部分的骨架。

\begin{remark}[符号动力—SFT—Markov 覆盖实例簇]\label{rem:typed-address-biaxial-completion-cluster-symbolic}
若把 \(\mathcal X\) 取为符号动力系统或其 sofic/Markov 覆盖，则时间推进天然由移位给出；同步、复位、解析深度与在线延迟可显式计算；\(q\)-重碰撞矩、pressure、\(\zeta\) 与 primitive 轨道谱则直接进入有限状态算子框架。
在本文语境里，\(\Fold_m\)、同步核、primitive 轨道谱与 \(\zeta_B\) 都属于这一簇的核心接口。
\end{remark}

\begin{remark}[逆极限—adic—solenoid 实例簇]\label{rem:typed-address-biaxial-completion-cluster-inverse-limit}
若把多尺度一致性真正做成一个逆系统
\[
X_\infty\cong \varprojlim X_m,
\]
则分辨率提升便可被严格解释为同一对象的更细近似，而不只是修辞。
定理 \ref{thm:recursive-addressing-dyadic-inverse-limit-stable-object} 已经给出逆极限地址与稳定点对象之间的严格对象化；在这一簇中，前缀一致性、restriction 相容性与层间自然性都具有真正的对象层语义。
这类实例最适合承载多尺度一致性、地址精化、账本前缀与紧群/全不连通纤维的双轴几何。
\end{remark}

\begin{remark}[覆盖—层—Čech 实例簇]\label{rem:typed-address-biaxial-completion-cluster-cech}
若把 \(\mathcal U\) 做成真正的覆盖与局部截面系统，则局部对象与全局对象的分野最清楚；\(\mathrm{NULL}\)、失败证书、重叠相容性与闭路 residual 也都可被精确定位。
这一路径的关键接口，正是定义 \ref{def:prefix-site-h2-gluing-obstruction}、命题 \ref{prop:null-as-h2-obstruction} 与上一节中的闭路 residual 语义。
凡是要把相位、干涉、规范差与局部—全局障碍写成硬数学，这一簇都不可绕过。
\end{remark}

\begin{remark}[层析—核化—Hilbert 接口实例簇]\label{rem:typed-address-biaxial-completion-cluster-hilbert}
若把局部扫描、边界剖面、有限秩核化与再现核空间结合起来，则局部读出可通过有限秩或低秩对象被审计，混合碰撞核、moment-kernel、Hankel--Prony 与 RKHS 特征嵌入也可进入统一空间。
这一簇尤其适合承载“粗糙可见量—表观随机—统计可区分性”的桥梁角色。
\end{remark}

\begin{remark}[代价半环—资源预序—函子化实例簇]\label{rem:typed-address-biaxial-completion-cluster-budget}
若把 \(\mathcal C\) 做成真正的资源函子或代价预序，则空间、时间、谱与账本四类预算可同时进入同一比较框架；
上一节关于 \(d_{\mathrm{res}}\) 的准伪度规性、不可对称性与不可替代性，也就获得真正的对象层宿主。
这正是“几何为何首先更像预算学而不是测地学”的最自然实例化方向。
\end{remark}

\subsubsection{相位—账本双轴模板}\label{subsubsec:typed-address-biaxial-completion-phase-ledger-template}
在本文当前主线里，最值得单独抽出来的是相位圆因子与 prime-ledger 因子的二分结构。
它不是附属现象，而是一个非常稳定的实例化模板。

\begin{proposition}[相位—账本双轴模板]\label{prop:typed-address-biaxial-completion-phase-ledger-template}
设某个实例具有：
\begin{enumerate}
  \item 一个连通的相位因子 \(K_{\mathrm{phase}}\)；
  \item 一个全不连通的外置账本因子 \(K_{\mathrm{ledger}}\)；
  \item 一个短正合扩张
  \[
  0\to K_{\mathrm{ledger}}\to E\to K_{\mathrm{phase}}\to 0;
  \]
  \item 层间相容的前缀读出与逆极限一致性。
\end{enumerate}
则它自然提供一类适合本文的双轴实例：相位负责连续读出、局部相容与 holonomy 语义；账本负责残差外置、历史可逆与 prime-register 方向。
特别地，在 prime-localized 情形下，这一模板局部退化为
\[
0\to \prod_{p\in S}\ZZ_p\to \widehat{G_S}\to \TT\to 0,
\]
即“连通圆相位 + 全不连通素数账本核”的二分结构。
\end{proposition}

\begin{proof}
定理 \ref{thm:killo-localization-circle-dimension-prime-ledger-splitting} 已经给出
\[
0\to \prod_{p\in S}\ZZ_p\to \widehat{G_S}\to \TT\to 0
\]
的精确分解，并说明连通圆因子与全不连通 prime-ledger 核在结构上不可混同。
因此，连通因子给出相位与连续参数的最小宿主，全不连通因子给出不可有限化的残差/账本方向，而短正合列则把二者组织成同一对象，却又不把它们误认成同一种几何量。
若再加上层间相容的前缀读出，这一扩张便自然承载本文需要的双轴结构。证毕。
\end{proof}

命题 \ref{prop:typed-address-biaxial-completion-phase-ledger-template} 的价值在于，它把 prime-register、圆维、solenoid 与相位 residual 这几条此前分散的线压成了一个简洁模板：
真正适合本文的对象，往往既不是纯相位对象，也不是纯账本对象，而是两者的扩张。

\subsubsection{逆极限一致性是硬门槛}\label{subsubsec:typed-address-biaxial-completion-inverse-limit-threshold}
任何多尺度叙事都喜欢说“不同分辨率看到的是同一对象”，但在本文里，这句话必须是严格命题，而不能只是比喻。

\begin{proposition}[多尺度一致性门槛]\label{prop:typed-address-biaxial-completion-multiscale-threshold}
若某个实例声称“分辨率 \(m\) 的差异只影响精细度，而不改变本体”，则必须存在一组限制映射
\[
\pi_{m',m}:X_{m'}\to X_m
\qquad (m'\ge m)
\]
使它们构成逆系统，并且存在对象
\[
X_\infty\cong \varprojlim X_m
\]
作为跨尺度一致本体。
若这一条件失败，则“不同分辨率看到的是同一对象”只是一句非严格口号，而不是对象层事实。
\end{proposition}

\begin{proof}
命题 \ref{prop:recursive-addressing-refinement} 已经说明：前缀细化若要保持在同一可见 \(\sigma\)-代数与同一对象层内，就必须服从 restriction 相容性。
定理 \ref{thm:recursive-addressing-dyadic-inverse-limit-stable-object} 则进一步说明：只有当这类限制映射真正组成逆系统时，跨层一致地址与稳定点对象才具备严格对象化语义。
因此，若没有 \(\pi_{m',m}\) 的相容系统与其逆极限对象，所谓“同一本体在不同尺度的读出”便无法被严格定义。证毕。
\end{proof}

这条门槛非常重要，因为它把“粗糙世界”“多尺度几何”“同一对象的不同分辨率”全部从直觉说法推回了数学约束。
没有逆系统，相位、账本、曲率候选与统计冻结都只能是彼此平行的散段结构，无法汇总成统一时空骨架。

\subsubsection{非例：哪些对象不够格}\label{subsubsec:typed-address-biaxial-completion-nonexamples}
为避免误用，这里顺手指出三类常见但不充分的对象。

\begin{remark}[纯集合编码不足]\label{rem:typed-address-biaxial-completion-nonexample-set}
若某个对象只给出状态到某个集合的单射编码，却没有更新、覆盖、残差与预算结构，那么它至多提供实现层地址，不能自动产生 \((\preceq,\tau,d_{\mathrm{res}},\Omega)\)。
\end{remark}

\begin{remark}[纯概率空间不足]\label{rem:typed-address-biaxial-completion-nonexample-probability}
若某个对象只给出一组随机变量与边缘分布，而没有局部—全局拼接、失败证书与多尺度一致性，那么它只能承载统计外观的一部分，不能承载本文的完整粗糙时空骨架。
\end{remark}

\begin{remark}[纯光滑流形不足]\label{rem:typed-address-biaxial-completion-nonexample-manifold}
一个预先给定的光滑流形当然可以承载经典几何，但若它没有 residual 外置、\(\mathrm{NULL}\) 分解、谱接口与资源代价，那么它描述的是一种更强、更规整、也更后验的几何。
它可以是本文骨架的某种极限收缩像，却不是当前阶段的最小实例模板。
\end{remark}

\subsubsection{首选最小实例簇}\label{subsubsec:typed-address-biaxial-completion-preferred-minimal-cluster}
结合全文现有结构，最自然的首选实例并不是单独某个对象，而是一个最小复合实例簇。

\begin{definition}[首选最小实例簇]\label{def:typed-address-biaxial-completion-preferred-minimal-cluster}
定义
\[
\mathcal I_{\min}
:=
\bigl(
\mathcal P_{\mathrm{LOST}},
\mathcal X_{\mathrm{inv}},
\mathcal U_{\mathrm{Cech}},
\mathcal R_{\mathrm{spec}},
\mathcal E_{\mathrm{prime}},
\mathcal C_{\mathrm{budget}}
\bigr),
\]
其中：
\begin{enumerate}
  \item \(\mathcal P_{\mathrm{LOST}}\) 是观察者—状态—forcing—更新系统；
  \item \(\mathcal X_{\mathrm{inv}}\) 是由 \(X_m\) 与 restriction 映射组成的逆系统及其 \(X_\infty\)；
  \item \(\mathcal U_{\mathrm{Cech}}\) 是局部覆盖、局部截面与重叠转移系统；
  \item \(\mathcal R_{\mathrm{spec}}\) 是由 \(q\)-碰撞矩、pressure、\(\zeta\) 与 primitive 轨道谱组成的谱接口；
  \item \(\mathcal E_{\mathrm{prime}}\) 是由纤维残差、\(\mathrm{NULL}\) 证书与动态 prime-register 组成的外置层；
  \item \(\mathcal C_{\mathrm{budget}}\) 是把空间/时间/谱/账本预算组织成资源预序的代价层。
\end{enumerate}
\end{definition}

定义 \ref{def:typed-address-biaxial-completion-preferred-minimal-cluster} 的目的并不是“一步推出全部物理学”，而是满足一个更克制但也更坚实的目标：
它已经足以同时承载局部性、因果性、多尺度一致性、粗糙统计、相位 residual 与资源几何。
换言之，若要在本文主线上寻找一个真正最小、但又不失真的物理型实例簇，这组数据几乎是不可再减的。

\subsubsection{小结：宿主选择规则}\label{subsubsec:typed-address-biaxial-completion-instantiation-summary}
本节真正完成的是进入物理解释之前的门禁设置。
其结论可以压缩为五点。

第一，并非任何数学对象都能实例化物理时空骨架；它必须同时保留 forcing、局部性、因果、失败结构、谱接口与资源预算。

第二，允许的实例化应被写成一个六元组
\[
(\mathcal P,\mathcal X,\mathcal U,\mathcal R,\mathcal E,\mathcal C),
\]
而不是单一流形、单一概率空间或单一编码。

第三，只要实例化满足观察者对称、因果相容、多尺度一致、局部—全局障碍、残差外置与谱—资源可审计六条条件，就可自然抽取出粗糙时空四元组
\[
(\preceq,\tau,d_{\mathrm{res}},\Omega).
\]

第四，最适合本文的具体对象，不是纯相位对象，也不是纯账本对象，而往往是“连通相位因子 + 全不连通 prime-ledger 因子”的扩张。

第五，经典几何若出现，只能是这些结构进一步收缩、对称化与可积化后的特殊情形，而不是实例化的起点。

这样，上一节给出的粗糙时空四元组终于获得了严格的宿主选择规则。
顺着这一点，下一节便可把这一模板真正落到一个首选最小实例簇上；在此之后，同步延迟、传播上界、dyadic holonomy 与证书闭环，都应被读成该实例簇上的具体实现。
